\section{The dihedral family}\label{the-dihedral-family}

Let \(n\ge3\), let \(V=\mathbb Z/n\mathbb Z\), and put
\(U_n=(\mathbb Z/n\mathbb Z)^\times\). For \(k\in U_n/\{\pm1\}\) let \[
E_k=\{\{x,x+k\}:x\in V\}.
\] Each \(E_k\) is a Hamiltonian cycle. If \(n\) is odd, all vertices
have one type. If \(n\) is even, even and odd vertices have separate
types. Each \(E_k\) is a further type. Vertex--edge incidence is
containment, edges of different types are mutually incident, and in the
even case every even vertex is incident to every odd vertex. These are
exactly the incidence systems of {[}L126, §§3.1--3.2{]}.

\textbf{Theorem R1.Dihedral.} For \(n>3\), the full correlation group is
the affine group \[
C(\Gamma_n)=V\rtimes U_n,
\qquad x\longmapsto ax+b.
\tag{D1}
\] Its type-preserving subgroup consists of the affine maps with
\(a=\pm1\), with arbitrary \(b\) for odd \(n\) and even \(b\) for even
\(n\). This gives the exact pair
\((\operatorname{Aut}D_{2n},\operatorname{Inn}D_{2n})\), where the
subscript is the group order. For \(n=3\) it gives the weak
representation stated in the source.

\textbf{Proof.} First the vertices are recognizable. For even \(n\),
their two types have size \(n/2\), while every edge type has size \(n\).
For odd \(n>3\), put \(d=\varphi(n)/2\ge2\). A vertex has degree \(2d\),
and an edge object has degree \(2+n(d-1)\). Their difference is
\((n-2)(d-1)>0\). Thus every correlation preserves \(V\) in both cases.
Each edge object is determined by its two endpoints.

The affine maps in (D1) preserve incidence and send \(E_k\) to
\(E_{ak}\). They act transitively on incident vertex--edge pairs: the
image of the pair \((0,\{0,1\})\) may be prescribed by choosing its
vertex image \(b\) and one of the two endpoint differences \(a\in U_n\).
There are \(n\varphi(n)\) such pairs. A correlation fixing one pair
fixes its edge type, hence acts on its Hamiltonian cycle fixing an edge
and both endpoints. An automorphism of a cycle fixing two adjacent
vertices is the identity, since successive neighbors are forced. It
fixes all vertices and therefore all edge objects. The action on these
pairs is thus free as well as transitive, proving (D1).

An affine map fixes every edge type only if it fixes \(E_1\),
equivalently \(a=\pm1\); these two multipliers indeed fix every \(E_k\).
For even \(n\) it fixes both vertex types exactly when \(b\) is even.
For odd \(n\) there is no further constraint.

Write \(D_{2n}=\langle r,s:r^n=s^2=1,\ srs=r^{-1}\rangle\). Its
automorphisms are exactly \[
r\mapsto r^a,\qquad s\mapsto r^b s,
\quad a\in U_n, b\in V.
\] Their composition is the affine composition law. Conjugation by \(r\)
has parameters \((1,2)\) and conjugation by \(s\) has parameters
\((-1,0)\), so the inner subgroup is precisely the subgroup just
computed. These formulas also verify compatibility of the two pair
isomorphisms.

For \(n=3\) the incidence graph is a hexagon. Its type-preserving group
is \(S_3\) and its full group has order twelve. Since
\(\operatorname{Aut}S_3=\operatorname{Inn}S_3=S_3\), the same inclusion
into the type-preserving subgroup gives compatible injections.
\(\square\)

\textbf{Source qualifications.} The unit group \(U_n\) need not be
cyclic. The source's factor \(C_{\varphi(n)}\) must be replaced by
\(U_n\); for instance \(U_8\cong C_2^2\). The printed equality involving
\(\operatorname{Out}D_6\) in the small-case proof is also unnecessary:
\(\operatorname{Out}S_3\) is trivial. Neither correction changes the
asserted pair in Theorem 3.1.

For some parameters \(\Gamma_n\) is only an incidence system. At
\(n=5\), two differently colored disjoint edges form a flag with no
vertex completing it. At \(n=6\), opposite-parity vertices at difference
three are incident but have no incident edge object in common. This does
not invalidate Theorem 3.1, whose defined notion of representation uses
incidence systems. If a genuine finite geometry is required, the
pair-preserving cleaning and chamber-completion constructions of Díaz
Ramos--Molinier--Viruel {[}F010{]} apply. Their complete proofs are included
in the common paper; no archived construction theorem is left unpaid.

\section{Symmetric groups and the exceptional
degree}\label{symmetric-groups-and-the-exceptional-degree}

For \(n\ge2\) let \(X_n\) be the rank-two incidence geometry of vertices
and edges of the complete graph on \(n\) vertices. The point and edge
classes have sizes \(n\) and \(\binom n2\).

\textbf{Theorem R1.Symmetric.} If \(n\ge4\), then \[
C(X_n)=K(X_n)=S_n.
\] Consequently \(X_n\) is an exact representation of \(S_n\) for
\(n\ge4\), \(n\ne6\). For \(n=3\) it is the preceding weak
representation. The literal \(n=2\) instance is not a representation of
\(S_2\) in the source's sense.

\textbf{Proof.} For \(n\ge4\) the two type sizes differ. Correlations
therefore preserve the vertices. Edges are uniquely recovered from their
endpoint sets, and all vertex permutations extend. This proves the
displayed pair. We retain the classical theorem
\(\operatorname{Aut}S_n=\operatorname{Inn}S_n\) for \(n\ge4\),
\(n\ne6\), also used in {[}L126, Theorem 4.1{]}. For \(n=2\), both
groups of \(X_2\) have order two, whereas both \(\operatorname{Inn}S_2\)
and \(\operatorname{Aut}S_2\) are trivial. Thus the source's parameter
exclusion \(n\ne3,6\) needs a small-parameter qualification. A
one-object rank-one geometry supplies the trivial pair when desired.
If the parameter $n=1$ is included, the declared rank-two model has an
empty edge type and hence fails the source's surjectivity requirement
for the type map; the same one-object rank-one replacement represents
the trivial group.
\(\square\)

We can make the exceptional construction exact, without assuming an
order calculation for \(\operatorname{Aut}S_6\). A \emph{duad} is a
two-element subset of a six-element set, and a \emph{syntheme} is a
partition into three duads. Form \(Q\) with the fifteen duads as points,
fifteen synthemes as lines, and containment as incidence, as in {[}L126,
pp.5--6{]}. Every point lies on a line and every line has three points,
so \(Q\) is a geometry.

\textbf{Theorem R1.Six.} The duad--syntheme geometry is an exact
representation of \(S_6\).

\textbf{Proof.} Distinct duads are collinear exactly when they are
disjoint. Hence a type-preserving automorphism preserves disjointness of
duads. The largest families of pairwise intersecting duads are precisely
the six stars consisting of all five duads through a fixed symbol.
Indeed, if an intersecting family has no common symbol, choose
\(\{a,b\}\) and \(\{a,c\}\) in it and then a member avoiding \(a\). That
member must be \(\{b,c\}\), and any further member must meet all three
edges of this triangle, so lies in the triangle. Such a family has at
most three members. The six five-member stars therefore recover the six
original symbols. Every type-preserving automorphism is the
corresponding permutation of symbols; synthemes are then determined by
their incident duads. Thus \(K(Q)=S_6\).

In the abstract group \(S_6\), the involution classes of transpositions,
products of two disjoint transpositions, and products of three have
sizes fifteen, forty-five, and fifteen respectively. This follows by
choosing the moved symbols and pairing them. Therefore every
automorphism of \(S_6\) permutes the first and third classes as a pair.
A transposition commutes with a triple transposition exactly when it is
one of its three factors: a commuting transposition must exchange two
symbols paired with each other by the triple transposition. Consequently
the duad--syntheme incidence relation is intrinsically the commuting
relation between these two classes. Every automorphism of \(S_6\) acts
on \(Q\) as a correlation, and this action is faithful because
transpositions generate \(S_6\).

Conversely every correlation of \(Q\) normalizes \(K(Q)\) and hence acts
by conjugation on its identified \(S_6\). This conjugation action is
faithful. A type-preserving element in its kernel belongs to the trivial
center of \(S_6\). A type-swapping element in its kernel would induce an
\(S_6\)-equivariant bijection between duads and synthemes. Their
stabilizers cannot be conjugate: a duad stabilizer contains seven
transpositions, whereas a syntheme stabilizer contains only its three
constituent transpositions. Thus no such bijection exists.

The two embeddings are inverse because the first action sends each inner
automorphism to the natural symbol permutation, and conjugation by its
image acts on that subgroup as the original group automorphism. Hence
the pair is exactly \((\operatorname{Aut}S_6,\operatorname{Inn}S_6)\).
\(\square\)

The native finite replay also gives \(|C(Q)|=1440\) and \(|K(Q)|=720\).
These counts corroborate the proof; they are not needed to deduce the
pair isomorphism. We do not rederive the atlas examples for \(A_6\) or
\(A_7\) cited but not constructed in {[}L126{]}. General finite
existence is supplied by common Theorem UR.FinitePairs, the theorem of
D\'iaz Ramos--Molinier--Viruel [F010]. For a finite group $L$, take
$G=\Aut(L)$ and $H=\operatorname{Inn}(L)$, which is normal in $G$.
The \(A_5\) case appears next with a small explicit construction.

\section{Cube and octahedron}\label{cube-and-octahedron}

Regard the cube graph as the Hamming graph on \(\{0,1\}^3\). Give even
and odd vertices distinct types, and give edges and square faces two
further types. First take the source's literal incidence induced from
the cube: vertex--edge, vertex--face and edge--face incidence are
containment, and different vertex objects are not incident.

\textbf{Theorem R1.Cube.} The literal rank-four incidence system has
pair \[
(C,K)\cong(S_4\times C_2,S_4),
\tag{C1}
\] with the second factor in \(C\) generated by the antipodal map. It
represents \(S_4\times C_2\) exactly. Adding incidence between
opposite-color vertices precisely when they are adjacent cube vertices
gives a geometry with the same pair and twenty-four chambers.

\textbf{Proof.} The type sizes are four, four, twelve and six, so every
correlation preserves the vertex union, the edges, and the faces. The
vertex--edge truncation recovers the cube graph. Edges and faces have
distinct vertex shadows, and every graph automorphism preserves their
families and the parity bipartition. Recognition therefore identifies
the full correlation group with the cube graph group.

Translations of \(\{0,1\}^3\) and coordinate permutations give a group
of order \(8\cdot6=48\). An automorphism fixing zero permutes its three
neighbors. Once these are fixed, each weight-two vertex is the common
neighbor of the appropriate pair other than zero, and the weight-three
vertex is forced. Thus the displayed group is the full graph group. The
parity-preserving subgroup has order twenty-four and acts faithfully on
the four even vertices, hence is \(S_4\). Translation by \((1,1,1)\) is
central, interchanges parity, and supplies the direct factor in (C1).

For \(L=S_4\times C_2\), its center is the second factor and is
characteristic. Every automorphism consequently has the form \[
(s,\epsilon)\longmapsto
(a s a^{-1},\epsilon\,\chi(s)),
\] where \(a\in S_4\) and \(\chi:S_4\to C_2\) is either trivial or the
sign character. The description follows by passing to
\(L/Z(L)\cong S_4\), using completeness of \(S_4\), and then inspecting
the central coordinate. The sign shear is a central involution in
\(\operatorname{Aut}L\), and the other factor is its inner group.
Identifying it with the antipodal map proves compatibility of (C1) with
the required pair.

The literal system has no chamber, since its two vertex types are never
incident. In the repaired system, each chamber consists of an edge, its
two endpoints and one of its two incident faces. There are
\(12\cdot2=24\). Every flag extends: within a square, two
opposite-parity vertices are adjacent; every vertex--edge or edge--face
flag extends through an endpoint or a face; and a vertex--face flag
extends through either boundary edge at that vertex. These observations
also handle singleton and empty flags. The added adjacency is already
recovered by the edge objects, so it changes neither group. \(\square\)

The repair establishes geometry but does not assert firmness: some
panels have a unique completion. The source's group-pair theorem
survives for its literal incidence system even though its use of the
word geometry does not.

\section{\texorpdfstring{Dodecahedron, icosahedron and
\(A_5\)}{Dodecahedron, icosahedron and A\_5}}\label{dodecahedron-icosahedron-and-a_5}

Use the Petersen graph with vertices the two-element subsets of a
five-element set and adjacency given by disjointness. It has ten
vertices and fifteen edges. Its twelve pentagons correspond to
Hamiltonian cycles of the complete graph on the five symbols: the five
edges of such a Hamiltonian cycle are Petersen vertices, and
disjointness puts them in pentagon order. Conversely the five pairs in a
Petersen pentagon use each symbol twice, since pairs through one symbol
form an independent set of size at most two in a pentagon. They
consequently form a Hamiltonian cycle on the five symbols.

The natural \(S_5\) acts transitively on these twelve cycles. A cycle
stabilizer is a dihedral group of order ten, all of whose permutations
on the five symbols are even. Thus \(A_5\) has two orbits \(\mathcal F\)
and \(\mathcal P\), each of size six, and every odd permutation
interchanges them. These are the face and Petrie families in the
source's hemidodecahedral construction.

Take vertex, edge, \(\mathcal F\) and \(\mathcal P\) as four types. Use
containment for their vertex and edge incidences. For a face
\(F\in\mathcal F\) and a Petrie cycle \(P\in\mathcal P\), set \[
F*P\quad\Longleftrightarrow\quad F\text{ and }P\text{ have a common Petersen edge}.
\tag{P1}
\] The rule (P1) is an explicit completion of an incidence left
unspecified in {[}L126, p.7{]}; it is not quoted as a printed rule.

\textbf{Theorem R1.Petersen.} The completed system is a geometry with
120 chambers and group pair \((S_5,A_5)\). It represents both \(A_5\)
and \(A_5\times C_2\) exactly.

\textbf{Proof.} For completeness, the required small cycle calculation
has two outcomes:

\begin{center}
\begin{tabular}{rrr}
\toprule
Pairs in $\mathcal F\times\mathcal P$ & Common vertices & Common edges\\
\midrule
30&3&2\\6&0&0\\\bottomrule
\end{tabular}
\end{center}

Here is the complete six-case calculation with one cycle fixed as
$(1,2,3,4,5)$. We write $ab$ for the Petersen vertex $\{a,b\}$.
Rotations and reversal identify Hamiltonian cycles; the six opposite-family
cycles and their common Petersen vertices are
\[
\begin{array}{c|c@{\qquad}c|c}
12354&12,23,45&12435&12,15,34\\
12543&12,34,45&13245&15,23,45\\
13524&\varnothing&14325&15,23,34.
\end{array}
\]
Each listed ordering has odd sign; the twelve Hamiltonian cycles split
into two six-element orbits because their dihedral stabilizers consist
of even permutations. Thus the list is complete. In each three-vertex
intersection, disjointness gives a path of length two. Transitivity on
the first family gives the table for all thirty-six pairs. The optional
replay enumerates all twelve cycles and verifies the same calculation.

Every Petersen edge belongs to two cycles in each family, by
transitivity and the count \(6\cdot5/15=2\). Each incident face--Petrie
pair has a common edge, and each of their common vertices lies on a
common edge by the table. Hence every flag extends to a
vertex--edge--face--Petrie chamber. The number of chambers is \[
15\cdot2\cdot2\cdot2=120.
\]

The type sizes ten, fifteen, six and six recognize vertices and edges.
The graph group is \(S_5\): its largest independent sets are the five
four-member stars of pairs through a symbol. As in the duad proof, a
pairwise intersecting family with no common symbol has size at most
three. Recovering these stars recovers the five symbols. Faces and
Petrie objects have distinct vertex supports, so restriction is
injective. All of \(S_5\) preserves (P1), while exactly \(A_5\) fixes
both cycle families. The recognition lemma gives the asserted group
pair.

Finally \(\operatorname{Aut}A_5=S_5\) and \(\operatorname{Inn}A_5=A_5\),
a classical retained group fact. For \(L=A_5\times C_2\), the center is
characteristic and the quotient is \(A_5\). There is no nontrivial
homomorphism \(A_5\to C_2\), since \(A_5\) is perfect. Thus
\(\operatorname{Aut}L=\operatorname{Aut}A_5=S_5\), with inner subgroup
\(A_5\). The same pair proves both representations. \(\square\)

Among \(S_5\)-invariant completions of the source's two specified
truncations, (P1) is the only one that gives a geometry. The thirty and
six cross-pairs are the two orbits. To see transitivity on the thirty,
fix one Hamiltonian cycle. Its order-five rotation is transitive on the
five noncomplementary opposite-family cycles: the only Hamiltonian edge
sets fixed by this rotation are that cycle and its complement. Including
the six disjoint pairs produces flags having no common point, while
excluding the thirty intersecting pairs leaves no chambers. This
uniqueness concerns that precise completion problem, not all possible
representations of these groups.
